A formulação robusta transportada da Matéria é expressa por
\begin{equation}
 \min_{\mathbf{x}\in\mathcal D}
 \left(\overline V(\mathbf{x}),\overline U_{fin}(\mathbf{x})\right),
 \qquad \overline g_j(\mathbf{x})\leq0,\quad j=1,\ldots,6,
\label{eq:problema_robusto}
\end{equation}
em que $\mathcal D$ reúne os limites das cinco variáveis e a admissibilidade geométrica. Para cada candidata, calculam-se
\begin{equation}
 \widetilde x_i^{(c)}=x_i(1+\rho\xi_i^{(c)}),\quad
 \xi_i^{(c)}\sim\mathcal U(-1,1),\qquad
 \overline f_m=\frac{1}{N_c}\sum_{c=1}^{N_c}f_m(\widetilde{\mathbf{x}}^{(c)}),\quad
 \overline g_j=\frac{1}{N_c}\sum_{c=1}^{N_c}g_j(\widetilde{\mathbf{x}}^{(c)}).
\label{eq:agregacao_robusta}
\end{equation}
Na implementação atual, a matriz de multiplicadores é reutilizada durante a otimização. O protocolo de comparação requer matrizes idênticas em R e C e sementes distintas entre repetições. Como a acomodação geométrica é recalculada, uma perturbação pode alterar o número inteiro de peças: essa formulação representa sensibilidade do processo de dimensionamento e não apenas tolerâncias de fabricação de uma ponte com número de peças fixo. A verificação posterior de uma solução construída deve preservar sua configuração discreta.

O NSGA-II combina seleção por viabilidade, ordenação por não dominância e preservação de diversidade por distância de aglomeração \parencite{Deb2002}. A distância é a soma dos intervalos entre vizinhos normalizados em cada objetivo, com ordenação própria por objetivo; não é uma medida de hipervolume. Os descendentes são avaliados antes da seleção elitista na união das populações, conforme o Algoritmo~\ref{alg:comparacao}.

\begin{table}[!ht]
\centering
\caption{Configuração inicial transportada da Matéria para o piloto; o orçamento final será definido pela verificação de convergência.}\label{tab:otimizacao_piloto}
\begin{tabular}{lc}
\toprule
Parâmetro & Configuração inicial \\
\midrule
População / gerações & 50 / 150 \\
Inicialização & Aleatória uniforme \\
Cruzamento SBX & $p_c=0{,}90$, $\eta_c=15$ \\
Mutação polinomial & $\eta_m=20$ \\
Eliminação de duplicatas & Sim \\
Realizações geométricas, $N_c$ & 30 \\
Amplitude relativa, $\rho$ & 5\% \\
Repetições propostas por cenário & 10 sementes pareadas (1 a 10) \\
\bottomrule
\end{tabular}
\end{table}

Esses parâmetros são um ponto de partida, não evidência de suficiência do orçamento. A probabilidade efetiva de mutação, a versão das bibliotecas, os limites dimensionais e as tolerâncias serão registrados na configuração executável. O piloto comparará pelo menos dois orçamentos de busca; se diferenças entre materiais forem da ordem da dispersão entre sementes, será necessário ampliar a busca antes de interpretá-las.

\begin{algorithm}[!ht]
\caption{Comparação pareada e reavaliação estrutural}\label{alg:comparacao}
\begin{algorithmic}[1]
\Require Base auditada, sistema de classes, ações, vãos, domínio e sementes
\For{cada espécie, vão e semente}
  \State Fixar propriedades complementares e matriz de perturbações comum
  \For{cenário $S\in\{R,C\}$}
    \State Definir a rigidez $E_S$ e inicializar e avaliar a população $P_0$
    \For{cada geração}
      \State Gerar descendentes $Q_t$ por seleção, SBX e mutação
      \State Avaliar acomodação, cargas, esforços, objetivos e restrições de $Q_t$
      \State Agregar as avaliações conforme a Equação~\ref{eq:agregacao_robusta}
      \State Selecionar em $P_t\cup Q_t$ por viabilidade, frente e aglomeração
    \EndFor
    \State Reavaliar soluções finais nominalmente e com amostra independente
    \State Selecionar menor volume nominal entre soluções que atendem ao critério comum
  \EndFor
  \State Reavaliar a geometria $\mathbf{x}_C$ com rigidez $E_R$, sem reotimizar
  \State Registrar $\Delta V$, $U_{fin}^{C\rightarrow R}$, $U_Q^{C\rightarrow R}$ e restrições
\EndFor
\State Agregar diferenças pareadas entre sementes e entre espécies
\end{algorithmic}
\end{algorithm}

O volume nominal da Equação~\ref{eq:volume_modelo} será usado em $\Delta V$, enquanto o objetivo de busca robusta é $\overline V$; ambos serão armazenados para evitar misturar estimativas. Uma solução aceita pelas restrições médias deverá ser novamente verificada em sua geometria nominal. Se nenhuma solução atender ao critério comum, o caso será identificado como sem solução viável encontrada, sem interpretar isso como prova de inexistência de projeto viável.

A verificação independente da otimização compreende: comparação de respostas com soluções analíticas em geometria fixa; conferência do equilíbrio das cargas distribuídas entre longarinas; avaliação das envoltórias nos limites dos trechos de carregamento; repetição entre sementes; e reavaliação das soluções selecionadas com perturbações não utilizadas na busca. O critério de aceitação da amostra de verificação deverá ser fixado antes da campanha e distinguido das restrições médias do otimizador.
